\subsection{重新考察膜}\label{sec:11j}
固定 $(V,\allowbreak I_V)$ 上的\emph{二次复体积形式}，即复直线之间的同构
\[
\eta_V^2:\lambda_{\CC}^{\mathrm{top}}(V)^{\otimes2}\longrightarrow\CC.
\]
等价地，可将它视为在相差符号的意义下确定的复体积形式 $\pm\eta_V:\lambda_{\CC}^{\mathrm{top}}(V)\to\CC$。相关的\emph{平方相位映射}——在特殊 Lagrangian 几何中广为人知——为
\begin{equation}\label{eq:11.34}
\alpha_V:\operatorname{Gr}(V)\to S^1,\qquad
\alpha_V(\Lambda)=\frac{\eta_V(v_1\wedge\cdots\wedge v_n)^2}{|\eta_V(v_1\wedge\cdots\wedge v_n)|^2},
\end{equation}
其中 $v_1,\allowbreak \ldots,\allowbreak v_n$ 是 $\Lambda$ 的任意一组基。该定义良好，因为非零元素 $v_1\wedge\cdots\wedge v_n\in\lambda_{\CC}^{\mathrm{top}}(V)$ 在相差 $\R^\times$ 倍的意义下与基无关，而平方去除了符号歧义。\znote{170}{平方相位映射是特殊 Lagrangian 几何的基本对象。nLab 关于 Lagrangian Grassmann 流形的条目指出，不稳定地 $H^1(\operatorname{LGrass}(\CC^n),\Z)\cong[U(n)/O(n),U(1)]$，由 $\det^2:U(n)\to U(1)$ 给出——因为 $O(n)$ 的行列式为 $\pm1$，$\det^2$ 过商良好。$\eta_V(v_1\wedge\cdots\wedge v_n)^2$ 正是 $\det^2$ 的具体实现，所以\eref{11.34}的 $\alpha_V$ 表示 Maslov 类 $\mu$；这也解释了为何取平方恰好消去基的符号歧义。从分次的角度看，$\alpha_V$ 的相位正是膜结构中实值相位 $\alpha^\#$ 的底：$\exp(2\pi i\alpha^\#)=\alpha_V(\Lambda)$ 把 $\alpha_V$ 拉回到 $\R$ 上成为多值函数，其单值化所需的整值修正由 Maslov 指标控制，所以 $\alpha_V$ 同时含有 $\Z/2$ 的符号信息与 $\Z$ 的 Maslov 信息。[来源：nLab "Lagrangian Grassmannian"，https:/\allowbreak /\allowbreak ncatlab.org/\allowbreak nlab/\allowbreak show/\allowbreak Lagrangian+\allowbreak Grassmannian（2026-11-15 访问）]}
\begin{definition*}
\emph{线性 Lagrangian 膜}是三元组 $(\Lambda,\allowbreak \alpha^\#,\allowbreak P^\#)$，由 Lagrangian 子空间 $\Lambda$、满足 $\exp(2\pi i\alpha^\#)=\alpha_V(\Lambda)$ 的实数 $\alpha^\#$，以及 $\operatorname{Pin}_n$ 主齐性空间 $P^\#$ 连同同构 $P^\#\times_{\operatorname{Pin}_n}\R^n\cong\Lambda$ 组成。称 $\alpha^\#$ 为\emph{实值相位}或\emph{分次}。
\end{definition*}
还需要考虑这种膜的族。具体地，设 $E\to B$ 为满足 $2c_1(E)=0$ 的辛向量丛。该条件保证，一旦固定相容复结构 $I_E$，复线丛 $\lambda_{\CC}^{\mathrm{top}}(E)^{\otimes2}$ 就
\SourcePage{178}
平凡，从而 $E$ 容许逐纤维的二次复体积形式。固定一个这样的形式 $\eta_E^2$，以 $\alpha_E$ 表示 Lagrangian Grassmann 流形丛 $\operatorname{Gr}(E)\to B$ 上相应的平方相位映射。
\begin{definition*}
Lagrangian 子丛 $F\subset E$ 的\emph{膜结构}，由光滑函数 $\alpha^\#:B\to\R$ 和 $F$ 上的 Pin 结构 $P^\#$ 组成，其中 $\exp(2\pi i\alpha^\#(b))=\alpha_E(F_b)$。
\end{definition*}
显然，膜结构存在的障碍是由平方相位 $b\mapsto\alpha_E(F_b)$ 表示的类 $\mu_F\in H^1(B)$，以及 $w_2(F)\in H^2(B;\Z/2)$。若这些障碍消失，则 $F$ 上膜结构的同构类构成 $H^0(B)\times H^1(B;\Z/2)$ 上的仿射空间。对应运算是：给分次加上一个局部常值的 $\Z$ 值函数，以及将 Pin 结构与实线丛作张量积。回到初始数据，应指出 $E$ 的子丛上的膜结构概念依赖于 $\eta_E^2$。更确切地说，只取决于它作为 $\lambda_{\CC}^{\mathrm{top}}(E)^{\otimes2}$ 平凡化的同伦类：变形 $\eta_E^2$ 时，平方相位函数连续变化，并可提升为实值相位的同伦。因此，实际的选择自由度由 $H^1(B)$ 给出。
\begin{note}\label{st:11.18}
令 $\xi_E\subset\lambda_{\CC}^{\mathrm{top}}(E)$ 为使 $\pm\eta_E$ 取实值的一维实子空间。$\xi_E$ 的非平凡性是为 $\eta_E^2$ 找到全局平方根的障碍。对任意 $b\in B$ 及 $F_b$ 的任意基 $\{v_k\}$，有 $v_1\wedge\cdots\wedge v_n\in\alpha_E(F_b)^{1/2}\xi_E$。假定给定膜结构，将它乘以 $\exp(-\pi i\alpha^\#)$，便得到同构 $\lambda^{\mathrm{top}}(F)\cong\xi_E$。因此，分次确定 $F$ 相对于 $\xi_E$ 的一个定向。给分次加上 $\sigma\in\Z$，使该定向改变 $(-1)^\sigma$ 倍。
\end{note}

建立一般框架之后，回到较受限的情形：平凡向量丛 $E=B\times V$，且 $I_E=I_V$、$\eta_E^2=\eta_V^2$ 均为常值。在此背景下，现在的膜概念与先前定义的抽象膜有何关系？容易看出，映射\eref{11.34}表示 Maslov 类 $\mu_V$。此外，短正合列 $\Z/2\to\operatorname{Pin}_n\to O_n$ 诱导纤维化 $B\operatorname{Pin}_n\to BO_n$，纤维为 $B\Z/2=\RP^\infty$；它可与万有纤维化沿表示 $w_2$ 的映射 $BO_n\to K(\Z/2,\allowbreak 2)$ 的拉回识别。这只是 $w_2$ 与 Pin 结构关系的同伦理论表述。因此，可将 $\operatorname{Gr}(V)^\#$ 取为两个拉回 $\alpha_V^*(\R\to S^1)$ 与 $i_V^*(B\operatorname{Pin}_n\to BO_n)$ 的纤维积，其中 $i_V$ 是将 $\operatorname{Gr}(V)$ 嵌入通常的 $n$ 维子空间 Grassmann 流形，再嵌入 $BO_n$ 得到的映射。于是，对给定 $F\subset B\times V$，相应映射 $B\to\operatorname{Gr}(V)$ 到 $\operatorname{Gr}(V)^\#$ 的提升的同伦类，与二元组 $(\alpha^\#,\allowbreak P^\#)$ 的同构类之间存在双射。这在同伦层面回答了我们的问题，虽大体令人满意，但并不完全。不过，与其寻找形式上更完善的陈述，不如简要回顾主要应用；这会具体说明，膜结构的新定义在实际作用上等价于旧定义。
\SourcePage{179}
取一对线性 Lagrangian 膜 $\Lambda_k^\#=(\Lambda_k,\allowbreak \alpha_k^\#,\allowbreak P_k^\#)$（$k=0,\allowbreak 1$），使 $(\Lambda_0,\allowbreak \Lambda_1)$ 横截。如\secref{11h}，可赋予它们绝对指标和定向空间。具体地，取路径 $\lambda\in\Omega^-(\operatorname{Gr}(V),\allowbreak \Lambda_0,\allowbreak \Lambda_1)$，在相应丛 $F_\lambda\subset[0;1]\times V$ 上取膜结构 $(\alpha_\lambda^\#,\allowbreak P_\lambda^\#)$，使分次满足 $\alpha_\lambda^\#(k)=\alpha_k^\#$，并取同构 $f_k:(P_\lambda^\#)_k\cong P_k^\#$，两处均为 $k=0,\allowbreak 1$。其余完全如\eref{11.25}。值得细想其中一些细节，以澄清膜结构在哪里起作用。假设给定路径 $\lambda$。存在唯一函数 $\alpha_\lambda^\#:[0;1]\to\R$，满足 $\alpha_\lambda^\#(0)=\alpha_0^\#$ 和 $\exp(2\pi i\alpha_\lambda^\#(s))=\alpha_V(\lambda(s))$。按我们的规定，还须有 $\alpha_\lambda^\#(1)=\alpha_1^\#$；这对 $\lambda$ 相对端点的同伦类施加限制，实际上选出一个优先同伦类。相反，为 $F_\lambda$ 赋予 Pin 结构没有障碍；但考虑路径的单参数族时，在 $[0;1]^2$ 上的相应向量丛上相对边界赋予 Pin 结构，就会对同伦类施加非平凡限制，保证定向空间不会出现非平凡单值变换。
\begin{note}\label{st:11.19}
实际上，只取决于 $(P_\lambda^\#,\allowbreak f_0,\allowbreak f_1)$ 的同构类，所以有效地只有两种可能。要从一种变到另一种，只需将 $f_0$ 或 $f_1$ 乘以非平凡元素 $-e\in\ker(\operatorname{Pin}_n\to O_n)$；这将诱导定向空间的反定向自同构。
\end{note}
\begin{example}\label{st:11.20}
考察最低维情形 $n=1$。考虑带标准生成元 $v=1$ 的 $\R$。依定义，$\operatorname{Pin}_1=\{\pm e,\allowbreak \pm v\}\subset\operatorname{Cl}(\R)=\R[v]/(v^2-1)$，所以投影 $\operatorname{Pin}_1\to O_1$ 分裂；此后固定像为 $\{e,\allowbreak v\}$ 的分裂。因此，实线丛 $F\to B$ 上的 $\operatorname{Pin}_1$ 结构等价于在底空间 $B$ 上选择另一个实线丛。特别地，线性 Lagrangian 膜由实直线 $\Lambda\subset V$、实值相位 $\alpha^\#$ 和另一个抽象实直线 $\beta$ 组成。对于横截的两个这样的膜，
\begin{equation}\label{eq:11.35}
\begin{aligned}
i(\Lambda_0^\#,\Lambda_1^\#)&=\lfloor\alpha_1^\#-\alpha_0^\#\rfloor+1,\\
o(\Lambda_0^\#,\Lambda_1^\#)&\cong\Hom(\beta_0,\beta_1)\otimes\Lambda_1^{\otimes i(\Lambda_0^\#,\Lambda_1^\#)}.
\end{aligned}
\end{equation}
这里 $\lfloor\cdot\rfloor$ 为向下取整。只需选择合适路径，并展开\eref{11.25}，即可得到这些公式。
\end{example}

\subsection{平移运算}\label{sec:11k}
对任意 $\sigma\in\Z$，线性 Lagrangian 膜有一个 $\sigma$ 重平移运算，将 $\Lambda^\#=(\Lambda,\allowbreak \alpha^\#,\allowbreak P^\#)$ 映为
\begin{equation}\label{eq:11.36}
S^\sigma\Lambda^\#=(\Lambda,\alpha^\#-\sigma,P^\#\otimes\lambda^{\mathrm{top}}(\Lambda)^{\otimes\sigma}).
\end{equation}
$\sigma=1$ 时简记为 $S\Lambda^\#$。\znote{172}{平移 $S^\sigma$ 把实值相位减去 $\sigma$，同时把 Pin 结构张量上 $\lambda^{\mathrm{top}}(\Lambda)$ 的 $\sigma$ 次幂；两者必须同步进行，否则 $\exp(2\pi i\alpha^\#)=\alpha_V(\Lambda)$ 不再成立。这个运算说明分次并非一个绝对数字：它只在相差一个整数的意义下由 $\Lambda$ 决定，而 $S^\sigma$ 恰把这个选择整体平移。绝对指标 $i$ 正是把这种“相差整数”的分次固定成 $\Z$ 分级的工具，\eref{11.37}的第一条说明平移 $\sigma$ 使指标减少 $\sigma$。nLab 的 Maslov 指标条目指出该指标取整值，并对路径的同伦不变。[来源：nLab "Maslov index"，https:/\allowbreak /\allowbreak ncatlab.org/\allowbreak nlab/\allowbreak show/\allowbreak Maslov+\allowbreak index（2026-11-15 访问）]}其关键性质是：
\SourcePage{180}
\begin{lemma}\label{st:11.21}
对满足 $\Lambda_0\pitchfork\Lambda_1$ 的任意线性 Lagrangian 膜 $\Lambda_0^\#,\allowbreak \Lambda_1^\#$，
\begin{equation}\label{eq:11.37}
\begin{aligned}
i(\Lambda_0^\#,S^\sigma\Lambda_1^\#)&=i(\Lambda_0^\#,\Lambda_1^\#)-\sigma,\\
o(\Lambda_0^\#,S^\sigma\Lambda_1^\#)&\cong o(\Lambda_0^\#,\Lambda_1^\#).
\end{aligned}
\end{equation}
\end{lemma}
\begin{proof}
指标结论直接可得，故只考虑定向空间。为简单起见，假设 $\sigma=1$。在 Lagrangian Grassmann 流形中，取从 $\Lambda_0$ 到 $\Lambda_1$ 的路径 $\lambda$，以及从 $\Lambda_1$ 回到自身的环路 $\rho$，将其串接记为 $\lambda\cdot\rho$。相应地，在上半平面和闭圆盘上有算子 $D_{H,\allowbreak \lambda}$、$D_{H,\allowbreak \lambda\cdot\rho}$ 和 $D_{D,\allowbreak \rho}$，关系如下：作边界连通和，将 $s\in\partial H$、$s\gg0$ 与 $1\in\partial D$ 识别，则按此方式粘合 $D_{H,\allowbreak \lambda}$ 与 $D_{D,\allowbreak \rho}$ 所得算子可与 $D_{H,\allowbreak \lambda\cdot\rho}$ 识别。

现在假设 $\lambda$ 使函数 $\alpha_V(\lambda(t))$ 有一个实值提升，在端点取值 $\alpha_0^\#,\allowbreak \alpha_1^\#$。再假设 $\mu(\rho)=-1$，于是将该实值提升延拓到 $\lambda\cdot\rho$ 后，新端点的值为 $\alpha_1^\#-1$。接着考虑相应丛 $F_\lambda\to[0;1]$ 和 $\widetilde F_\rho\to[0;1]$；后者的记号与\secref{11i}一致。可以为它们赋予 Pin 结构 $P_\lambda^\#,\allowbreak \widetilde P_\rho^\#$ 和同构
\begin{equation}\label{eq:11.38}
\begin{cases}(P_\lambda^\#)_0\cong P_0^\#,\\(P_\lambda^\#)_1\cong P_1^\#,
\end{cases}\qquad
\begin{cases}(\widetilde P_\rho^\#)_0\cong P_1^\#,\\(\widetilde P_\rho^\#)_1\cong P_1^\#\otimes\lambda^{\mathrm{top}}(\Lambda_1).
\end{cases}
\end{equation}
利用\eref{11.38}的第二、第三个同构将两者结合，得到 $F_{\lambda\cdot\rho}$ 上的 Pin 结构，且在端点分别有到 $P_0^\#$ 和 $P_1^\#\otimes\lambda^{\mathrm{top}}(\Lambda_1)$ 的典范同构。

由\stref{引理}{11.11}及定向空间的定义，$F_\lambda$ 和 $F_{\lambda\cdot\rho}$ 上 Pin 结构的这些选择确定同构 $\det(D_{H,\allowbreak \lambda})\cong o(\Lambda_0^\#,\allowbreak \Lambda_1^\#)$ 与 $\det(D_{H,\allowbreak \lambda\cdot\rho})\cong o(\Lambda_0^\#,\allowbreak S\Lambda_1^\#)$。另一方面，$\widetilde P_\rho^\#$ 是\eref{11.32}意义下的扭曲 Pin 结构，由\stref{引理}{11.17}给出同构 $\det(D_{D,\allowbreak \rho})\cong\lambda^{\mathrm{top}}(\Lambda_1)$。结合\eref{11.11}，得到
\[
\begin{aligned}
o(\Lambda_0^\#,S\Lambda_1^\#)\otimes\lambda^{\mathrm{top}}(\Lambda_1)
&\cong\det(D_{H,\lambda\cdot\rho})\otimes\lambda^{\mathrm{top}}(\Lambda_1)\\
&\cong\det(D_{H,\lambda})\otimes\det(D_{D,\rho})\\
&\cong o(\Lambda_0^\#,\Lambda_1^\#)\otimes\lambda^{\mathrm{top}}(\Lambda_1).
\end{aligned}
\]
\end{proof}
\begin{note}\label{st:11.22}
也可不引用\stref{引理}{11.17}，而直接照用其部分证明。具体地，取某个分解 $\Lambda_1=\Lambda_{1,\allowbreak +}\oplus\Lambda_{1,\allowbreak -}$，其中 $\dim(\Lambda_{1,\allowbreak -})=1$，并令 $\rho(s)=\Lambda_{1,\allowbreak +}\oplus\exp(-\pi sI_V)\Lambda_{1,\allowbreak -}$。由\eref{11.20}可见，$(\delta_H)_{\lambda\cdot\rho}$ 比 $(\delta_H)_\lambda$ 多出一个张量因子 $\Lambda_{1,\allowbreak -}$。但扭曲 Pin 结构的选择恰对应于同一空间的定向，因而可识别 $(\delta_H)_{\lambda\cdot\rho}\cong(\delta_H)_\lambda$。此论证避免了指标理论，较为初等；不过，对于理解\eref{11.37}与\eref{11.26}中同构的相互作用等目的，原证明更方便。
\end{note}

\SourcePage{181}
\subsection{定向算子}\label{sec:11l}
将\stref{命题}{11.13}适配到膜的新定义很直接。具体地，仍取前述 $(S,\allowbreak \nu)$，并为 Lagrangian 子丛 $F_\nu\subset\partial S\times V$ 赋予膜结构。对每个条带状端 $\zeta$，便得到一对线性 Lagrangian 膜 $(\Lambda_{\zeta,\allowbreak 0}^\#,\allowbreak \Lambda_{\zeta,\allowbreak 1}^\#)$，其绝对指标和定向空间按\eref{11.26}确定 $D_{S,\allowbreak \nu}$ 的指标与行列式线。本节最后的任务是将这一结论推广到稍大的一类算子，以便用于非齐次伪全纯映射的线性化算子。

回顾\secref{8h}中定义非退化 $\bar\partial$ 算子的数据：$S$ 是带条带状端的 Riemann 曲面，$(E,\allowbreak \omega_E,\allowbreak J_E)\to S$ 是 Hermite 向量丛，$\nabla$ 是 $E$ 上的联络，$F\subset E|\partial S$ 是 Lagrangian 子丛。它们在条带状端上满足适当的渐近收敛条件。为方便应用，在每个端添加区间 $\epsilon_\zeta(\{\pm\infty\}\times[0;1])$ 来紧化，得到带角曲面 $\bar S$。以 $\overline{\partial S}$ 表示 $\partial S$ 在 $\bar S$ 中的闭包，即 $\partial S$ 与各点 $\epsilon_\zeta(\pm\infty,\allowbreak k)$（$k=0,\allowbreak 1$）的并。由渐近行为，上述结构分别在 $\bar S$ 上、最后一个则在 $\overline{\partial S}$ 上，有典范延拓 $(\bar E,\allowbreak \omega_{\bar E},\allowbreak J_{\bar E},\allowbreak \bar\nabla,\allowbreak \bar F)$；限制到无穷远区间时，这些延拓恰是最初所说的极限数据 $(E_\zeta,\allowbreak \omega_{E_\zeta},\allowbreak J_{E_\zeta},\allowbreak \nabla_\zeta,\allowbreak \Lambda_{\zeta,\allowbreak k})$。现在再添加以下数据：
\begin{itemize}
\item 另一个与 $\omega_{\bar E}$ 相容的复结构 $I_{\bar E}$，连同逐纤维二次复体积形式 $\eta_{\bar E}^2$ 及相应平方相位函数 $\alpha_{\bar E}$；
\item 子丛 $\bar F\subset\bar E|\overline{\partial S}$ 上的膜结构，由分次 $\bar\alpha^\#$ 和 Pin 结构 $\bar P^\#$ 组成。
\end{itemize}
取一个条带状端 $\zeta$。考察 $k=0,\allowbreak 1$ 对应的点 $\epsilon_\zeta(\pm\infty,\allowbreak k)$，上述数据使每个 $\Lambda_{\zeta,\allowbreak k}$ 成为线性 Lagrangian 膜 $\Lambda_{\zeta,\allowbreak k}^\#$。起初它们位于不同的辛向量空间 $E_{\zeta,\allowbreak k}$ 中，但可用辛联络 $\nabla_\zeta$ 的平行移动，结合相应纤维 $\bar E_z$、$z\in\epsilon_\zeta(\{\pm\infty\}\times[0;1])$ 上的平方相位映射 $\alpha_{\bar E}$，将 $\Lambda_{\zeta,\allowbreak 0}^\#$ 移到 $E_{\zeta,\allowbreak 1}$ 中，从而定义绝对指标和定向空间。我们断言：考虑到这一点后，算子 $\bar\partial_\nabla$ 满足\eref{11.26}的对应公式。

证明相当直接。选择辛平凡化 $\bar E\cong\bar S\times V$，使之在无穷远区间上与联络 $\nabla_\zeta$ 相容。接着，要将 $I_{\bar E}$ 变形成常值复结构 $I_V$，并相应地将 $\eta_{\bar E}^2$ 变形成 $\eta_V^2$。后一部分在 $H^1(S)$ 中有一个显然的障碍，但只要适当选取最初平凡化的同伦类，该障碍就消失。二次复体积形式的变形诱导平方相位映射的同伦，分次也应相应地随之移动。结果是，现在由映射 $\bar\rho:\overline{\partial S}\to\operatorname{Gr}(V)$ 给出的 Lagrangian 子丛 $\bar F$，作为
\SourcePage{182}
$S\times(V,\allowbreak \omega_V,\allowbreak I_V)$ 的子丛带有膜结构。最后，将 $J_{\bar E}$ 变形成常值复结构 $I_V$，将 $\bar\nabla$ 变形成平凡联络。这诱导 Fredholm 算子空间内从 $\bar\partial_\nabla$ 到标准 Dolbeault 算子 $\bar\partial$ 的同伦，将问题归结到先前考虑的特殊情形。$S=H$ 时，得到\stref{引理}{11.11}的一个小推广：
\begin{equation}\label{eq:11.39}
\operatorname{index}(\bar\partial_\nabla)=i(\Lambda_0^\#,\Lambda_1^\#),\qquad
\det(\bar\partial_\nabla)\cong o(\Lambda_0^\#,\Lambda_1^\#).
\end{equation}
\begin{definition*}
这一情形足够重要，值得单独命名：称 $\bar\partial_\nabla$ 为线性 Lagrangian 膜对 $(\Lambda_0^\#,\allowbreak \Lambda_1^\#)$ 的\emph{定向算子}。\znote{174}{定向算子是带 Lagrangian 边界条件的非退化 $\bar\partial$ 算子的一般名称：它的指标等于两个边界膜的绝对指标，它的行列式线与定向空间同构。nLab 关于 Maslov 指标的条目列举了这类指标定理的经典来源（Robbin–Salamon 的两篇文章），其中路径情形正是 $S=H$ 的特例。名称中的“定向”指\eref{11.39}第二条：算子的行列式线被膜数据定向化以后，Floer 理论中相交点所携带的符号才能一致地确定。[来源：nLab "Maslov index"，https:/\allowbreak /\allowbreak ncatlab.org/\allowbreak nlab/\allowbreak show/\allowbreak Maslov+\allowbreak index（2026-11-15 访问）]}
\end{definition*}
